% Einheit 3 – Klammern auflösen (bank/terme/e3.jsonl, zusammenbau v1.2)
\einheitenkopf*[e3][Klammern auflösen]{Klammern auflösen}
\setcounter{aufgabe}{12}


% test: terme-e3-k2-s8-v2
\lbtest{Kannst du das schon? Dann weiter zu „Ausklammern“}{
\lbtt{a)}{$5 \cdot (a + 3) - 2 \cdot (a - 1)$: \lbfeld}
}{a) $3a + 17$}


% erkennung: terme-e3-k1-s0-v1, terme-e3-k1-s0-v2, terme-e3-k1-s0-v3, terme-e3-k1-s0-v4
\begin{kbaufgabe}{Erkennen, was vor einer Klammer steht}
\begin{kbblock}
\kbanweisung{Was steht im Term … direkt vor der Klammer?}
\item[]\hspace*{-\leftmargin}\lbzelle{2}{\kbteil{a)}{$7 + (a - 2)$: \lbfeld}{}}\hfill\lbzelle{2}{\kbteil{b)}{$10 - (b + 4)$: \lbfeld}{}}\par
\end{kbblock}
\begin{kbblock}
\item[]\hspace*{-\leftmargin}\lbzelle{2}{\kbteil{c)}{$5 \cdot (x - 1)$: \lbfeld}{}}\hfill\lbzelle{2}{\kbteil{d)}{$2y - 6 \cdot (y + 1)$: \lbfeld}{}}\par
\end{kbblock}
\end{kbaufgabe}


% leiter: terme-e3-k2-s1-v1, terme-e3-k2-s1-v2, terme-e3-k2-s2-v1, terme-e3-k2-s3-v1, terme-e3-k2-s4-v1, terme-e3-k2-s5-v1, terme-e3-k2-s6-v1, terme-e3-k2-s7-v1, terme-e3-k2-s8-v1
\begin{kbaufgabe}{Klammern auflösen}
\begin{kbblock}
\item[]\hspace*{-\leftmargin}\lbzelle{2}{\kbteil{a)}{$5a + (2b - 1) =$ \lbfeld}{}}\hfill\lbzelle{2}{\kbteil{b)}{$5a + (2b - 3) =$ \lbfeld}{}}\par
\end{kbblock}
\begin{kbblock}
\kbanweisung{Löse die Klammer mit der Tabelle auf.}
\kbteil{c)}{$5 \cdot (x + 3)$}{}
\kbzweispaltig[0.34]{\sachtabelle{ccc}{$\cdot$ & $x$ & $+3$}{$5$ & \leerzelle & \leerzelle}}{\lbantwortrechts[0.4cm]{= \kbfeld}}
\end{kbblock}
\begin{kbblock}
\kbteil{d)}{Rechne $30 - (12 + 8)$ auf zwei Wegen: erst die Klammer ausrechnen und dann abziehen; dann Glied für Glied abziehen. Vergleiche beide Ergebnisse.}{}
\item[]\kbraum{2}
\end{kbblock}
\begin{kbblock}
\item[]\hspace*{-\leftmargin}\begin{lbflucht}\makebox[1.6em][r]{e)\hspace{0.2em}} & $5a - (2b + 4)$ & $=$ & \lbfeld \\[\lbfluchtabstand] \makebox[1.6em][r]{f)\hspace{0.2em}} & $7a - (2b - 5 + c)$ & $=$ & \lbfeld \\[\lbfluchtabstand] \makebox[1.6em][r]{g)\hspace{0.2em}} & $-2 \cdot (x + 5)$ & $=$ & \lbfeld\end{lbflucht}\par
\end{kbblock}
\begin{kbblock}
\kbanweisung{Löse die Klammer auf und fasse zusammen.}
\item[]\hspace*{-\leftmargin}\begin{lbflucht}\makebox[1.6em][r]{h)\hspace{0.2em}} & $5x + 2 \cdot (x + 6)$ & $=$ & \lbfeld\end{lbflucht}\par
\end{kbblock}
\begin{kbblock}
\kbanweisung{Löse die Klammern auf und fasse zusammen.}
\item[]\hspace*{-\leftmargin}\begin{lbflucht}\makebox[1.6em][r]{i)\hspace{0.2em}} & $3 \cdot (2x + 5) - 2 \cdot (x - 4)$ & $=$ & \lbfeld\end{lbflucht}\par
\end{kbblock}
\end{kbaufgabe}


% pflicht: terme-e3-k3-s1-v3, terme-e3-k3-s2-v3, terme-e3-k3-s3-v1, terme-e3-k3-s4-v1
\begin{lbkasten}
\begin{kbaufgabe}{Verstanden?}
\begin{kbblock}
\item[]\lbhinweis{Gemischt – erkenne selbst, was zu tun ist.}
\kbteil{a)}{Nils hat vier Klammern aufgelöst. Welche Ergebnisse können nicht stimmen? Begründe, ohne genau zu rechnen.}{}
\kbfrage{(1) $7 - (x + 2) = 5 - x$ \\ (2) $3 \cdot (y - 4) = 3y - 4$ \\ (3) $-(a - 6) = -a - 6$ \\ (4) $9 + (b - 1) = 8 + b$}
\item[]\kbraum{2}
\end{kbblock}
\begin{kbblock}
\kbteil{b)}{Entscheide bei jeder Aussage, ob sie wahr oder falsch ist. Begründe, ohne genau zu rechnen.}{}
\kbfrage{(1) Steht ein Plus vor der Klammer, kann man die Klammer immer weglassen. \\ (2) Steht ein Minus vor der Klammer, bleibt das Vorzeichen des ersten Glieds in der Klammer immer gleich. \\ (3) Eine Zahl vor der Klammer wird mit jedem Glied in der Klammer malgenommen.}
\item[]\kbraum{3}
\end{kbblock}
\begin{kbblock}
\kbteil{c)}{Die Figur zeigt ein Rechteck. Die Maße sind in cm angegeben. Schreibe einen Term für den Flächeninhalt mit Klammer. Schreibe den Term dann ohne Klammer.}{}
\kbzweispaltig[0.48]{\viereck[seiten={x+2,5,x+2,5}]{(0,0)}{(6,0)}{(6,2.5)}{(0,2.5)}}{\lbantwortrechts[1.5cm]{\kbfeld}}
\end{kbblock}
\begin{kbblock}
\kbteil{d)}{Ein Kinoticket kostet 9 €. Eine Klasse mit 24 Schülern nimmt noch $x$ Begleiter mit. Schreibe die Gesamtkosten in Euro als Term mit Klammer. Löse die Klammer auf.}{}
\item[]\kbraum{2}
\end{kbblock}
\end{kbaufgabe}
\end{lbkasten}
